% vim: ai sw=2 spell spelllang=en

\section{Building and Linking the Kernel}
\frame{\sectionpage}

\iffalse
\begin{frame}[fragile]{Compiling One Object}

  \heading{\cmd{make V=1 fs/ext4/dir.o} (shortened)}
  \begin{lstlisting}[language=bash]
gcc -Wp,-MMD,fs/ext4/.dir.o.d -nostdinc -I./arch/x86/include -I./include \
  -include ./include/linux/compiler-version.h -include ./include/linux/kconfig.h \
  -include ./include/linux/compiler_types.h -D__KERNEL__ -std=gnu11 \
  -fno-strict-aliasing -mno-sse -mno-red-zone -mcmodel=kernel -O2 -g ... \
  -DKBUILD_MODFILE='"fs/ext4/ext4"' -DKBUILD_BASENAME='"dir"' \
  -DKBUILD_MODNAME='"ext4"' -DMODULE -c -o fs/ext4/dir.o fs/ext4/dir.c
  \end{lstlisting}

  \pause

  \begin{itemize}[<+->]
    \item \cmd{-nostdinc} -- no libc, no \file{/usr/include}
    \item \file{kconfig.h} includes \file{autoconf.h} -- every file sees all
      \file{CONFIG_*}
    \item \cmd{-mno-sse} -- no FPU in the kernel
    \item \cmd{-MMD} -- dependency file (later)
    \item \cmd{-DMODULE} -- built as \cmd{m}
  \end{itemize}

\end{frame}
\fi

\begin{frame}{From Objects to \file{vmlinux}}

  \begin{itemize}[<+->]
    \item \file{dir/built-in.a} -- per directory
      \begin{itemize}[<+->]
	\item \emph{Thin} archive: references \file{.o} files, no copies
	\item Historically \file{built-in.o} (\cmd{ld -r}), until 4.13
      \end{itemize}
    \item \file{vmlinux.a} -- all top-level \file{built-in.a} together
    \item \file{vmlinux.o} -- \cmd{ld -r} of everything
      \begin{itemize}[<+->]
	\item \cmd{objtool} -- validates stack frames, generates ORC unwind
	  data, \dots
	\item \cmd{modpost} -- section mismatches (\file{.init} vs.\ the rest)
      \end{itemize}
    \item \file{scripts/link-vmlinux.sh}
      \begin{itemize}[<.->]
	\item Links vmlinux \emph{several times}. Why?
      \end{itemize}
  \end{itemize}

\end{frame}

\begin{frame}{Why Link \file{vmlinux} Several Times?}

  \begin{itemize}[<+->]
    \item \emph{kallsyms}: the kernel contains its own symbol table
      \begin{itemize}[<.->]
	\item Stack traces with function names, \file{/proc/kallsyms}
      \end{itemize}
    \item Chicken \& egg
      \begin{itemize}[<.->]
	\item Symbol table needs final addresses
	\item Adding the table moves the addresses
      \end{itemize}
    \item The solution
      \begin{enumerate}[<.->]
	\setcounter{enumi}{-1}
	\item Empty table
	\item \file{.tmp_vmlinux1} $\rightarrow$ table of the right \emph{size}
	\item \file{.tmp_vmlinux2} $\rightarrow$ table with the right \emph{addresses}
	\item \file{.tmp_vmlinux3} -- only if something still moved
	\item Final link with the correct table
      \end{enumerate}
    %\item BTF is generated from \file{.tmp_vmlinux1} too (later)
    \item \file{System.map} must match the embedded table (checked)
  \end{itemize}

\end{frame}

\begin{frame}{Linking the Kernel}{The Whole Picture}

  \begin{center}
  \begin{tikzpicture}[node distance=5mm and 7mm]
    \path[use as bounding box] (-1.2, -3.8) rectangle (12.5, 1.1);
    \node[bnode] (c) {*.c};
    \node[bnode, below=of c] (o) {*.o};
    \node[bnode, below=of o] (bi) {built-in.a};
    \node[bnode, below=of bi] (va) {vmlinux.a};

    \draw[barrow] (c) -- node[btool, right] {gcc -c} (o);
    \draw[barrow] (o) -- node[btool, right] {ar T} (bi);
    \draw[barrow] (bi) -- node[btool, right] {ar T} (va);

    \node<2->[bnode, right=18mm of va] (vo) {vmlinux.o};
    \draw<2->[barrow] (va) -- node[btool, above] {ld -r} node[btool, below] {objtool} (vo);

    \node<3->[bnode, above=of vo] (t1) {.tmp\_vmlinux1};
    \node<3->[bnode, above=of t1] (t2) {.tmp\_vmlinux2};
    \node<3->[bnode, above=of t2, dashed] (t3) {.tmp\_vmlinux3};
    \draw<3->[barrow] (vo) -- node[btool, right] {ld} (t1);
    \draw<3->[barrow] (t1) -- node[btool, right] {kallsyms} (t2);
    \draw<3->[barrow, dashed] (t2) -- node[btool, right] {kallsyms} (t3);

    \node<4->[bnode, right=20mm of t1, fill=yellow!20] (btf) {.btf.o};
    \draw<4->[barrow] (t1) -- node[btool, above] {pahole -J} (btf);

    \node<5->[bnode, right=20mm of t3] (vu) {vmlinux.unstripped};
    \draw<5->[barrow] (t3) -- node[btool, above] {ld} (vu);
    \draw<5->[barrow] (btf) -- (vu);
    \node<5->[bnode, right=10mm of vu, fill=green!20] (vm) {vmlinux};
    \draw<5->[barrow] (vu) -- node[btool, above] {"strip"} (vm);
    \node<5->[bnode, below=of vm] (sm) {System.map};
    \draw<5->[barrow] (vu) |- (sm);
  \end{tikzpicture}
  \end{center}

  \onslide<6->{\footnotesize \file{vmlinux} is a regular ELF executable.
    But it cannot be booted directly on \xes.}

\end{frame}

\begin{frame}{From \file{vmlinux} to \file{bzImage}}{\xes}

  \begin{columns}[c]
    \begin{column}{.45\textwidth}
      \begin{tikzpicture}[remember picture,node distance=4mm]
	\path[use as bounding box] (-2.4, -6.2) rectangle (5, .4);
	\node[bnode, fill=green!20] (vm) {vmlinux};
	\node<2->[bnode, below=of vm] (vb) {compressed/vmlinux.bin};
	\node<3->[bnode, below=of vb] (gz) {compressed/vmlinux.bin.zst};
	\node<4->[bnode, below=of gz] (pig) {compressed/piggy.o};
	\node<5->[bnode, below=of pig] (cv) {compressed/vmlinux};
	\node<6->[bnode, below=of cv] (bvb) {boot/vmlinux.bin};
	\node<7->[bnode, right=10mm of bvb] (setup) {boot/setup.bin};
	\node<8->[bnode, below=of bvb, fill=red!20] (bz) {boot/bzImage};

	\draw<2->[barrow] (vm) -- node[btool, right] {objcopy -O binary} (vb);
	\draw<3->[barrow] (vb) -- node[btool, right] {zstd / gzip / xz / \dots} (gz);
	\draw<4->[barrow] (gz) -- node[btool, right] {mkpiggy + as} (pig);
	\draw<5->[barrow] (pig) -- node[btool, right] {ld + decompressor} (cv);
	\draw<6->[barrow] (cv) -- node[btool, right] {objcopy} (bvb);
	\draw<8->[barrow] (bvb) -- node[btool, right] {tools/build} (bz);
	\draw<8->[barrow] (setup) |- (bz);
      \end{tikzpicture}
    \end{column}

    \begin{column}{.5\textwidth}
      \begin{itemize}
	\item<2-> \file{arch/x86/boot/}
	\item<3-> \tikz[remember picture,baseline=-.5ex]{\node[coordinate,xshift=-1.4em](RAW){};}Raw
	  binary, compressed
	\item<5-> \tikz[remember picture,baseline=-.5ex]{\node[coordinate,xshift=-1.4em](DEC){};}Wrapped
	  by a \emph{decompressor}
	  \begin{itemize}[<.->]
	    \item Own small kernel: decompresses and relocates (KASLR)
	  \end{itemize}
	\item<7-> \tikz[remember picture,baseline=-.5ex]{\node[coordinate,xshift=-1.4em](SUP){};}\file{setup.bin}
	  -- 16-bit real-mode code
	  \begin{itemize}[<.->]
	    \item Header for bootloaders (boot protocol)
	    \item Also an EFI PE image (EFI stub)
	  \end{itemize}
	\item<9-> \cmd{make install} copies \file{bzImage} to \file{/boot/vmlinuz-*}
      \end{itemize}
    \end{column}
  \end{columns}

  \begin{tikzpicture}[remember picture,overlay]
    \draw<3->[kwarrow] (RAW.west) to (gz.east);
    \draw<5->[kwarrow] (DEC.west) to (cv.east);
    \draw<7->[kwarrow] (SUP.west) to[relative,in=230,out=330] (setup.north);
  \end{tikzpicture}

\end{frame}

\begin{frame}[fragile]{The Result}

  \heading{\file{vmlinux} vs.\ \file{bzImage}}
  \begin{itemize}[<+->]
    \item \cmd{ls -l vmlinux arch/x86/boot/bzImage}
    \item \cmd{file vmlinux arch/x86/boot/bzImage}
  \end{itemize}

  \uncover<.(1)->{\heading{Booting the Result}}
  \begin{itemize}[<+->]
    \item No need to install, no need to reboot
    \item QEMU loads \file{bzImage} directly
      \begin{lstlisting}[language={}]
qemu-kvm -m 2G \
    -kernel arch/x86/boot/bzImage \
    -append 'root=/dev/vda' \
    -drive file=root.img,if=virtio
      \end{lstlisting}
  \end{itemize}

\end{frame}
